% Answers to Chapter 3's exercises, lectures/L03/appendix/b_exercises.md. The book's own answers:
% the lecture publishes none.
\answerchapter{3}{Interrupts}

\answer{c3:ex:1}
Polling makes the main loop ask every input, on every pass, whether anything has happened. Each
input added makes every pass longer, the CPU spends its time checking inputs that have not changed,
and the reaction time to any one input is the time of a whole pass: while the loop is busy
elsewhere, in a blocking delay say, a press is noticed late, and a pulse shorter than a pass can be
missed altogether. An interrupt moves the watching into hardware: the pin logic detects the event
and runs its ISR within a few clock cycles whatever \code{main()} is doing, so an input that is
idle costs nothing, and the ISR only has to flag the event for \code{main()} to handle.

\answer{c3:ex:2}
\code{EICRA} configures \emph{what} triggers each external interrupt: its sense bits
\code{ISCn1:ISCn0} choose low level, any change, falling edge or rising edge, for INT0 and INT1
separately. \code{EIMSK} configures \emph{whether} each one is enabled at all: its \code{INT0} and
\code{INT1} bits switch the two interrupts on. Neither register does the other's job, and neither
replaces the global \code{I} bit in \code{SREG} that \code{sei()} sets.

INT0 can be enabled without touching \code{EICRA}, but it will not do what you want.
\code{EICRA} is 0 after a reset, so \code{ISC01:ISC00} is \code{00}: a low-level interrupt,
requested for as long as the pin is held low. With the book's button wiring, D2 idles low (the
external pull-down) and goes high only while the button is held, so INT0 fires continuously while
the button is \emph{not} pressed: it is requested again the moment the ISR returns, and the CPU
executes just one instruction of \code{main()} between two runs of it. In the toggle program of
\secref{c3:a:3} the LED toggles hundreds of thousands of times a second and looks half-lit, and a
press stops it in whichever state it happens to be in.

\answer{c3:ex:4}
Nothing is broken: the edge is the wrong one for this circuit. The button is wired with an external
pull-down (\secref{c2:a:4}), so D2 reads low while the button is up and is pulled high while it is
held down. A press is therefore a rising edge, low to high, and a release a falling edge, high to
low. \code{ISC01:ISC00 = 10} selects the falling edge, so INT0 fires on the release, exactly as
the colleague sees. Rising edge (\code{11}) is right for this wiring; a falling edge would mean the
press only for a button wired the other way round, to ground against a pull-up. (Contact bounce can
add stray edges in both directions, but the clean falling edge is the release.)

\answer{c3:ex:5}
D4 is \code{PORTD4}: port D, bit 4. Port D's formula is \code{PCINTx = PORTDx + 16}:
\[
  x = 4 + 16 = 20,
\]
so D4 is \code{PCINT20}. It belongs to port D's group, \code{PCIE2}/\code{PCMSK2}, whose
interrupts all share \code{PCINT2_vect}. The three registers and bits to touch:
\begin{itemize}
  \item \code{PCICR}: set \code{PCIE2}, to enable pin change interrupts for port D's group.
  \item \code{PCMSK2}: set bit 4, \code{PCINT20}, to select D4 within the group;
    \code{(1U << PCINT20)} and \code{(1U << 4)} are the same bit.
  \item \code{SREG}: set the global \code{I} bit, with \code{sei()}.
\end{itemize}

\answer{c3:ex:6}
Keep a copy of the last \code{PINB} value in a \code{static} byte. On each interrupt, read
\code{PINB} once and XOR it with the copy: the bits that are 1 in the result are the pins that
changed since the last interrupt. AND-ing that with the new reading keeps the pins that changed
\emph{and} are now high, which with the book's wiring are the buttons just pressed; releases drop
out. Test the two button bits, D11 (\code{PORTB3}) and D12 (\code{PORTB4}), in the result, set a
\code{volatile} flag for each button found pressed, for \code{main()} to read and clear, and store
the new reading as the copy.

Only the button bits are tested, because \code{PINB} shows every pin of the port, the LED
outputs included. A copy that starts at 0 matches two buttons idling low; initializing it from
\code{PINB} in \code{setup()}, before \code{sei()}, also covers a button held down at start-up.
